% СПО ЛКС - TCP/IP
% Пынькин Д.А. (с) 2011
\documentclass[ignorenonframetext, hyperref={pdftex, unicode}]{beamer}

\usepackage[russian]{babel}
\usepackage[utf8]{inputenc}
\usepackage[T1]{fontenc}

\usepackage{beamerthemesplit}

%\usetheme{Pittsburgh}
%\usecolortheme{dolphin}

\usetheme{Copenhagen}
\usecolortheme{wolverine}


\usepackage{ulem}
%\usepackage{html}

\usepackage{verbatim}

\usepackage{ucs}
\usepackage{listings}
\lstloadlanguages{C, make, bash}

\lstset{escapechar=`,
	extendedchars=false,
	language=C, 
	tabsize=2, 
	columns=fullflexible, 
%	basicstyle=\scriptsize,
	keywordstyle=\color{blue}, 
	commentstyle=\itshape\color{brown},
%	identifierstyle=\ttfamily, 
	stringstyle=\mdseries\color{green}, 
	showstringspaces=false, 
	numbers=left, 
	numberstyle=\tiny, 
	breaklines=true, 
	inputencoding=utf8x,
	keepspaces=true,
	morekeywords={u\_short, u\_char, u\_long, in\_addr}
	}

\definecolor{darkgreen}{cmyk}{0.7, 0, 1, 0.5}

\lstdefinelanguage{diff}
{
    morekeywords={+, -},
    sensitive=false,
    morecomment=[l]{//},
    morecomment=[s]{/*}{*/},
    morecomment=[l][\color{darkgreen}]{+},
    morecomment=[l][\color{red}]{-},
    morestring=[b]",
}

\title[СПОЛКС (http://goo.gl/32cTB)]{Системное программное обеспечение\\локальных компьютерных сетей}
\author{Денис Пынькин}
\date{2011 -- 2012}
%\institution[БГУИР]{Белорусский государственный университет информатики и радиоэлектроники}
%\logo{\includegraphics[width=1cm]{logo-kafEVM.png}}


\subtitle[TCP/IP]{Стек протоколов TCP/IP\\(Продолжение)}

\begin{document}

\mode<all>{\begin{frame}
\titlepage
\begin{center}
e-mail: denis.pynkin@bsuir.by\\
\end{center}
\begin{center}
{\bfseries http://goo.gl/32cTB}

{\tiny СЧАСТЬЕ ДЛЯ ВСЕХ, ДАРОМ, И ПУСТЬ НИКТО НЕ УЙДЕТ ОБИЖЕННЫЙ!\\
(c)Стругацкие, Пикник на обочине}
\end{center}
\end{frame}
}

%
% Далее начинается сама презентация
%

\section{ARP}

\begin{frame}{ARP}
	\begin{block}{ARP (Address Resolution Protocol —- протокол определения адреса)}
		использующийся в компьютерных сетях протокол низкого уровня,  предназначенный для определения адреса канального уровня по известному адресу сетевого уровня.
	\end{block}
\end{frame}


\begin{frame}{ARP: протокол определения адреса}
	Сетевой интерфейс имеет аппаратный адрес (48-битное значение для Ethernet или Token ring). 
Фреймы,  которыми обмениваются на аппаратном уровне,  должны адресоваться к корректному интерфейсу. 
Однако TCP/IP использует собственную схему адрессации: 32-битные IP адреса. 
Знание IP адреса хоста не позволяет ядру послать датаграмму этому хосту.\\
Драйвер Ethernet должен знать аппаратный адрес пункта назначения,  чтобы послать туда данные. В задачу ARP входит обеспечение динамического соответствия между 32-битными IP адресами и аппаратными адресами,  используемыми различными сетевыми технологиями.
\end{frame}

\begin{frame}{Принцип работы}
	\begin{enumerate}
		\item Узел,  которому нужно выполнить отображение IP-адреса на локальный адрес,  формирует ARP запрос,  вкладывает его в кадр протокола канального уровня,  указывая в нем известный IP-адрес,  и рассылает запрос широковещательно.
		\item Все узлы локальной сети получают ARP запрос и сравнивают указанный там IP-адрес с собственным.
		\item В случае их совпадения узел формирует ARP-ответ,  в котором указывает свой IP-адрес и свой локальный адрес и отправляет его уже направленно,  так как в ARP запросе отправитель указывает свой локальный адрес.
	\end{enumerate}
\end{frame}

\begin{frame}{ARP кэш}
	Эффективность функционирования ARP во многом зависит от ARP кэша (ARP cache),  который присутствует на каждом хосте.\\
	В кэше содержатся Internet адреса и соответствующие им аппаратные адреса.\\
	Стандартное время жизни каждой записи в кэше составляет 2 минуты с момента создания записи.
\end{frame}

\begin{frame}[fragile]{Пример просмотра ARP-кэша}
\scriptsize
	\begin{verbatim}
# arp -a
host1 (10.0.0.113) at 00:0f:fe:d4:e6:bb [ether] on eth0
host2 (10.0.0.15) at 78:e3:ff:94:63:58 [ether] on eth0
host3 (10.0.0.1) at 00:2f:29:35:03:7c [ether] on eth0
host4 (10.0.0.8) at 00:15:45:0e:2f:82 [ether] on eth0
	\end{verbatim}
\normalsize
\end{frame}

\mode<all>{\begin{frame}{}
\Huge
\begin{center}
	Спасибо за внимание!
	\bigskip
	Вопросы?
\end{center}
\normalsize
\end{frame}
}

\end{document}
%Конец файла
